\section{QAC Generation Pipeline}

Our pipeline turns multilingual chemistry patent text into
question--answer--context (QAC) records for retrieval evaluation. It produces
questions in five languages, designed for cross-lingual retrieval and kept close
to what a practitioner would actually search, in both technical and general
modes. Figure~\ref{fig:qac-generation-flow} summarizes the generation and
evaluation flow.

\begin{figure*}[t]
    \centering
    \includegraphics[width=0.98\textwidth]{multilingual_chemistry_retrieval_pipeline_EDITABLE.png}
    \caption{Overview of the QAC generation pipeline: source documents are
    converted into grounded candidate questions, scored, and exported as
    retrieval-ready records. Icons adapted from Font Awesome, licensed under CC BY 4.0. }
    \label{fig:qac-generation-flow}
\end{figure*}
\paragraph{Corpus and context construction.}
We build on two patent-derived corpora that contribute complementary kinds of
multilinguality. \textit{Google Patents} supplies scale and language breadth:
23{,}787 published applications, predominantly WIPO (WO, 15{,}048) and European
(EP, 4{,}610) filings, with a sizeable Mexican (MX, 3{,}952) share, spanning five
languages (English, French, Spanish, German, Chinese), from which we take the
title and abstract as the evidence block. On this side the non-English versions
are Google's \emph{machine} translations, and because Chinese coverage was zero
we machine-translated 400 documents to seed it using GPT 5.5. \textit{EPO} supplies clean
cross-lingual parallelism: 11{,}315 granted B1 specifications, each published as a
\emph{human}-translated English/German/French triple ($\approx$3{,}770 per
language) with professionally translated first claims, which we add to the
evidence block; we introduce no translations of our own on this side.


\paragraph{Candidate generation.}
For each source document we prompt \texttt{gpt-5-mini} to generate three
candidate questions, each grounded in a completely by the passage. (Table~\ref{tab:benchmark-stats} reports the
resulting counts). Two choices shape the query distribution. First, a question's
\emph{target language} is sampled per document relative to the languages that
document exists in: \emph{random-existing} picks a language the patent already
has a version in, so the query has a same-language answer;
\emph{random-missing} picks one it does \emph{not}, yielding a pure cross-lingual
``no-home'' query; and \emph{all} generates the question in every 5 languages
we targeted. Second, each question is written in one of two \emph{modes}:
\emph{technical} questions ask for a single concrete fact, a parameter,
material, outcome, method, or structural detail, answerable by the document, while \emph{general} questions ask about the document's problem, approach,
or application in deliberately paraphrased, low-overlap wording. The mode prompts
are written in the target language so questions read natively. The prompts (English version) are available in  Figures~\ref{fig:prompt-technical}
and~\ref{fig:prompt-general}.

\paragraph{Scoring, human audit.}
The three candidates per document are scored not by the generator but by a
separate verifier, \texttt{claude-sonnet-4.6}, which rates each candidate
on faithfulness and on mode-specific quality; we keep the highest-scoring
candidate per document.(verifier prompts in
Figures~\ref{fig:verifier-faith}--\ref{fig:verifier-general}). Because the
verifier is itself a model, we audited it against a 97-item human-annotated
sample using criteria that mirror the LLM rubric; three rounds of annotator
feedback further refined both the generation and the verifier prompts. On the
final sample the human mean is \textbf{8.33/10} (technical \textbf{8.72}, general
\textbf{7.91}) and \emph{no} item was rejected (none scored in the 0--3 band;
Figure~\ref{fig:human-scores}). We then export the standard retrieval
artifacts.

